% ==========================================
% TABEL VI.1a HASIL FUNCTIONALITY TESTING - MOBILE
% ==========================================
\begin{longtable}{|p{1.2cm}|p{4.5cm}|p{5cm}|p{2cm}|}
\caption{Hasil \textit{Functionality Testing} - Mobile (Pelari)} \label{tab:functionality-testing-mobile} \\

\hline
\textbf{ID} & \textbf{Skenario Pengujian} & \textbf{Hasil yang Diharapkan} & \textbf{Status} \\ \hline
\endfirsthead

\caption[]{Hasil \textit{Functionality Testing} - Mobile (Pelari) (lanjutan)} \\
\hline
\textbf{ID} & \textbf{Skenario Pengujian} & \textbf{Hasil yang Diharapkan} & \textbf{Status} \\ \hline
\endhead

\hline
\multicolumn{4}{r}{\textit{Bersambung ke halaman berikutnya}} \\
\endfoot

\hline
\endlastfoot

TC-01 & Pelari bergerak di sepanjang rute & ETA, jarak tersisa, waktu tempuh, dan \textit{pace} terhitung serta diperbarui secara dinamis & Berhasil \\ \hline

TC-02 & Pelari menekan tombol pusatkan kamera (\textit{center}) & Peta otomatis bergeser dan memusatkan tampilan ke titik lokasi pelari saat ini & Berhasil \\ \hline

TC-03 & Pelari menekan tombol \textit{emergency} (SOS) & Tampilan darurat muncul dan hitung mundur peringatan dimulai & Berhasil \\ \hline

TC-04 & Pelari menekan tombol menu utama & Laci menu terbuka dan menampilkan opsi navigasi tambahan & Berhasil \\ \hline

TC-05 & Pelari menekan tombol info tim & Panel informasi muncul menampilkan data tim dan anggota & Berhasil \\ \hline

TC-06 & Pelari mengubah tuas (\textit{switch}) \textit{visibility} & Status visibilitas berubah antara \textit{public} atau \textit{private} & Berhasil \\ \hline

TC-07 & Pelari menekan tombol \textit{logout} & Sesi pengguna berakhir dan aplikasi kembali ke layar \textit{login} & Berhasil \\ \hline

TC-08 & Pelari menekan layar (\textit{tap}) 2x saat kondisi \textit{emergency} & Kondisi \textit{emergency} dibatalkan (\textit{cancel}) dan kembali ke mode peta & Berhasil \\ \hline

TC-09 & Pelari menekan navigasi arah (sebelum/selanjutnya) & Peta berpindah dan memfokuskan (\textit{spectating}) pelari lain dalam tim yang sama & Berhasil \\ \hline

\end{longtable}